%=== CHAPTER ONE ===
%% Rewritten 2026-09-18 for the eventless first-moment interface.

\chapter{Introduction}
\label{ch:introduction}
\begin{spacing}{1.5}
\setlength{\parskip}{0.22in}

\section{Background}

Quadrotors are increasingly used for inspection, transport and delivery because
they can hover and manoeuvre in confined three-dimensional environments. Their
agility, however, is obtained from a light airframe with limited thrust and
torque margins. A package pickup can therefore change the vehicle dynamics on the
same time scale as the control loop. The total mass changes the
thrust-to-acceleration map, an offset load produces a persistent moment about the
composite centre of mass, and the payload lever arm changes the rotational
inertia.

Model predictive control is attractive for this problem because it predicts
future motion while handling actuator limits and coupled dynamics explicitly
\cite{rawlings2017mpc,mayne2014mpc}. For a quadrotor the relevant equations are
nonlinear and the attitude evolves on the unit-quaternion manifold; nonlinear MPC
therefore avoids the loss of fidelity caused by repeated linearisation over
aggressive motion \cite{grune2017nmpc,kamel2017linear}. Its performance
nevertheless depends on the model supplied to the prediction. If the controller
continues to use the empty-airframe parameters after grasping a parcel, even a
numerically converged optimisation solves the wrong problem.

This dissertation couples moving horizon estimation with nonlinear model
predictive control. MHE estimates states and parameters by optimising over a
recent window of measurements while enforcing the system dynamics and physical
bounds \cite{rao2003mhe}. Within that window the estimator identifies two
payload-dependent scalars, and from them the controller reconstructs every
payload-dependent term in its own model.

\section{Problem Statement}

Consider a quadrotor that follows a three-dimensional figure-eight path, grasps
an eccentric payload, transports it and releases it at a selected phase of the
trajectory. The controller has no access to the true payload parameters. Only
ordinary onboard flight signals are available: odometry, attitude, body rate and
rotor-speed telemetry.

Three model errors appear at attachment:
\begin{enumerate}
\item the composite mass changes abruptly, shifting the hover thrust and the
translational acceleration gain;
\item the payload lever arm increases roll and pitch inertia, reducing effective
inner-loop bandwidth;
\item the horizontal centre-of-mass offset creates a non-zero trim moment and
consumes part of the available torque authority.
\end{enumerate}
At release all three changes occur in the opposite direction.

The obvious way to handle release is to tell the controller when it happens. The
controller in this work in fact \emph{issues} the release command itself, so the
information is available for free. The central observation of this dissertation
is that using it is a mistake.

\begin{quote}
A release command is a \emph{request}. The model must describe the
\emph{vehicle}. The two diverge whenever a gripper is slow, jams, releases only
part of a load, or when a load is lost with no command at all --- and a
controller that clears its payload model on the command flies a loaded vehicle
with an empty model, silently, because nothing in its loop contradicts it.
\end{quote}

The central research question is therefore:

\begin{quote}
Can a quadrotor identify, from its ordinary flight signals alone, that a payload
has been grasped and that it has physically left --- with no attach or release
notification of any kind entering the estimator or the model update --- and
remain a constrained, well-tracking predictive controller throughout?
\end{quote}

This question turns out to have a specific technical obstacle, which
Chapter~\ref{ch:model} addresses and which shapes the whole design. The natural
payload parameterisation --- a payload mass and a grasp offset --- becomes
singular at exactly the moment of interest: as the payload mass goes to zero the
offset is undefined, and an estimator holding the offset fixed can only cancel
the resulting phantom trim moment by corrupting its mass estimate. The remedy
adopted here is to estimate the \emph{first mass moment} $\svec=\mP\bm r_{xy}$,
the product that the equations of motion actually contain, rather than its two
factors.

\section{Research Objectives}

The dissertation has six objectives:
\begin{enumerate}
\item derive a consistent 13-state nonlinear quadrotor model with explicit
payload mass, inertia and centre-of-mass entry points, and show which
parameterisation of those entry points remains well posed through release;
\item formulate a moving horizon estimator that augments the composite mass and
the horizontal first mass moment, deriving the inertia increment and the
centre-of-mass offset from that pair without any horizontal attachment geometry;
\item reconstruct the physical wrench from rotor speeds so that the estimator is
never driven by a signal computed from its own output;
\item formulate a predictive controller whose dynamics, equilibrium-input
baseline and inner-loop command scaling are all updated from one atomic,
health-gated estimate frame;
\item define a grasp--release lifecycle in which no attach or release
notification is admitted as evidence, and in which an unconfirmed release is held
as \emph{unresolved} rather than assumed complete;
\item evaluate the complete interface in software-in-the-loop against
command-triggered and command-armed baselines, under delayed detachment,
unsignaled loss, partial loss and thrust-map mismatch.
\end{enumerate}

\section{Scope}

The work covers simulation and software-in-the-loop validation of a rigidly
attached parcel. Airframe inertia, rotor geometry and the thrust and
yaw-reaction coefficients are treated as calibrated plant quantities. The
calibrated airframe mass appears in the rigid-body model, in the derived payload
mass $\mP=(\mT-\mB)^{+}$ and in offline scoring.

The information boundary is stated precisely, because a looser statement would be
false. \emph{No task-specific payload mass and no attach or release notification
enter the estimator or the controller's model update.} One mass-valued platform
constant does remain in the closed loop --- the airframe's \SI{0.30}{\kg} payload
envelope, used for an attach-time inertia floor and for a physical admissibility
bound --- and it is disclosed as a specification of the airframe rather than as
information about the parcel. Chapter~\ref{ch:model} states exactly where it
appears.

The scope does not include deformable or sloshing loads, cable-suspended payload
dynamics, wind, motor failure, hardware flight, or a formal closed-loop stability
proof for the estimator--controller interconnection. These are revisited in
Chapter~\ref{ch:conclusion}.

\section{Contributions}

The contributions are:
\begin{enumerate}
\item \textbf{A first-moment payload state.} Augmenting $\svec=\mP\bm r_{xy}$
rather than a mass-and-offset pair removes the release singularity, supplies the
centre-of-mass offset as $\cvec=\svec/\mT$ with no grasp geometry, and prevents
the estimator from encoding payload \emph{presence} into the mass channel.

\item \textbf{A well-posed reduction of the inertia model.} The second-order
horizontal terms of the parallel-axis increment are shown to carry a factor
$1/\mP$ in the first-moment coordinates, and are discarded for a quantified
\SI{2.6}{\percent} bias. The remaining mass-dependent coefficient is shown to act
as a lever by which the optimiser can dilute a phantom moment, and is frozen
within each window to break that derivative.

\item \textbf{A non-circular physical-input reconstruction}, including the
yaw-reaction channel, so that the estimator observes what the actuators did
rather than what the controller intended.

\item \textbf{A window vertical-force-balance re-initialisation} that uses only
the two velocity endpoints of the window and therefore requires no numerical
differentiation, reducing the seed error from \SI{17.1}{\percent} to
\SI{2.8}{\percent} in aggressive flight.

\item \textbf{An eventless payload lifecycle} in which the estimator owns
presence and may declare release, the controller owns the model and clears it
only on persistent confirmation, and an unconfirmed release becomes a hold with a
smooth brake to hover rather than an assumption of completion.

\item \textbf{A release rule built on a self-formed, frozen moment reference},
with mandatory moment evidence --- a requirement introduced after a
weaker rule caused real premature releases --- and a documented record of what
each threshold was selected against.

\item \textbf{A single atomic, health-gated payload interface} feeding both the
predictive model and a continuous body-rate command scaling with explicit
headroom limiting, replacing a design in which parameters arrived on separate
topics and the autopilot's own gains were rewritten in flight.

\item \textbf{A paired experimental evaluation} that separates release
\emph{semantics} from persistent confirmation, quantifies what eventlessness buys
(detection of unsignaled loss) and what it costs (confirmation delay, and a
severe sensitivity to thrust-map calibration), and reports the failure modes
rather than only the successes.
\end{enumerate}

\section{Dissertation Organisation}

Chapter~\ref{ch:literature} reviews predictive control, online estimation,
inertial-parameter identification and payload-adaptive aerial control.
Chapter~\ref{ch:model} derives the model and the payload parameter structure, and
is where the first-moment argument and its consequences are developed in full.
Chapter~\ref{ch:framework} presents the coupled architecture, the numerical form
of both optimal control problems, the continuous payload interface and the
eventless lifecycle. Chapter~\ref{ch:evaluation} describes the implementation and
reports the frozen experimental campaign, including the conditions under which
the interface fails. Chapter~\ref{ch:conclusion} summarises the findings and
states what remains open.

\end{spacing}
\newpage
